%======================================================================
\chapter{Data Collection and Integration}
%======================================================================
The project integrates multiple authoritative \gls{od} sources and implements rigorous processing techniques to develop a robust framework for spatial analysis and infrastructure planning. See Figure \ref{fig:data-pipeline} in Appendix for data pipeline visualization.

\section{Geographic and Infrastructure Data}

\subsection{Boundary Data}

%The boundary data was sourced from the \gls{row-open-data} \cite{region}. Geographic and infrastructure data formed the foundational layer of our research. Utilizing boundary information from the Open Data, we employed the \gls{utm-17n} to ensure precise geo-spatial measurements.

The foundation of our analysis rests on comprehensive geographic data from the \gls{row-open-data} \cite{rowdata2024}. We employed the \gls{utm-17n} system to ensure precise spatial measurements throughout our analysis. This geographic framework proved essential for accurate coverage calculations and infrastructure planning, particularly in areas where service boundaries intersect with major transportation corridors.

\subsection{Charging Infrastructure}

%The charging infrastructure mapping relied on the \href{https://openchargemap.org}{OpenChargeMap API} \cite{open}. It revealed a network of 183 charging stations within the \glsxtrshort{kwc-cma} bounds, with 169 being Level 2 chargers. Our data collection went far beyond simple station counting, capturing intricate details such as port configurations, operator information, and power specification metrics.

Our infrastructure assessment through the \gls{ocm} \cite{ocm2024} revealed a network of 183 charging stations across the KWC-CMA region. The existing infrastructure comprises 167 \gls{l2} stations providing 424 charging ports, operating at the standard 19.2 kW per port capacity specified by HoneyBadger Charging (2023) \cite{honey2023}. Additionally, 16 \gls{l3} stations offer 66 fast-charging ports, each capable of delivering up to 150 kW as documented by Future Energy (2022) \cite{future2022}. The spatial distribution of these stations formed a crucial baseline for our optimization model.

\subsection{Potential Locations}

% Potential charging locations were identified through an in-depth analysis using \gls{osm} via OSMnx\cite{boeing_2024_modeling}. The research targeted key areas including commercial centres, parking facilities, and public venues. This approach allowed for a sophisticated multi-factor location scoring system that assessed land use compatibility and evaluated grid proximity, creating a comprehensive view of available spaces within the \glsxtrshort{kwc-cma}.

\Gls{osm} data, accessed through OSMnx \cite{boeing2024}, provided the basis for identifying potential charging station locations. Our analysis considered commercial centres, retail locations, public parking facilities, and community venues, with particular attention to grid infrastructure proximity as detailed in \glsxtrshort{ieso} planning documents \cite{ieso2021}. The evaluation of each location incorporated land use compatibility, site accessibility, and property characteristics, creating a comprehensive dataset for optimization modelling.

\subsection{Transit Network Data}

% The transit network component of the research drew from Grand River Transit’s Open Data\cite{a2017_open}, providing a comprehensive capture of transportation infrastructure. This included a detailed mapping of bus routes and stop locations, ION Light Rail Transit infrastructure, service frequency metrics, and an extensive coverage area analysis. A sophisticated multi-modal accessibility scoring system was developed to integrate transit network considerations into the charging station optimization strategy.

% The entire data collection and processing workflow was meticulously managed through a dedicated \verb|src/data_manager.py| script, with critical location analysis performed in the\\ \verb|notebooks/02_location_analysis.ipynb| Jupyter notebook. Every stage of data collection adhered to rigorous geo-spatial standards, ensuring high-integrity, spatially aligned, and comprehensively analyzed datasets.

% By leveraging multiple data sources and employing advanced analytical techniques, the research created a robust methodology for identifying optimal \glsxtrshort{ev} charging station locations. This approach combines geographical, demographic, infrastructural, and transportation network insights to provide a comprehensive strategy for \glsxtrshort{ev} charging infrastructure development.

The integration of \gls{grt-open} \cite{rowdata2024} enabled sophisticated transit accessibility analysis. We examined bus route networks, stop locations, and ION Light Rail Transit infrastructure to develop a multi-modal accessibility scoring system. This integration proved crucial for understanding the relationship between public transit patterns and optimal charging infrastructure placement.

\section{Population and Electric Vehicle Data}

% Demographic analysis drew from Statistics Canada’s Census Profile\cite{a2021_profile} and \gls{row-open-data} \cite{region}, examining population characteristics at the census tract level. The investigation captured critical metrics including population density, housing characteristics, and demographic contextual factors. Geo-spatial integration ensured precise UTM projection alignment and accurate spatial joins, providing a nuanced understanding of the region's population dynamics.

% \Glsxtrshort{ev} ownership analysis utilized the Electric Vehicles in Ontario- By Forward Sortation Data\cite{a2024_electric} database which details the total number of \glsxtrshort{ev}s in Ontario. The data obtained was covering all \glsxtrshort{kwc-cma} \glsxtrshort{fsa} codes. The comprehensive processing included density calculations, Battery Electric Vehicle (BEV) and Plug-in Hybrid Electric Vehicle (PHEV) ratio analysis, spatial distribution mapping, and predictive growth trend projections. This approach allowed for a forward-looking assessment of \glsxtrshort{ev} adoption and potential charging infrastructure needs.

Our demographic analysis synthesized data from \gls{sc}'s 2021 Census \cite{statcan2021} with electric vehicle ownership patterns from \gls{odc}'s \glsxtrshort{fsa}-level database \cite{ontario2024}. The census data provided critical information about population distribution, tract boundaries, and housing characteristics across the region. This demographic foundation was enriched by detailed \glsxtrshort{ev} registration data, allowing us to map current ownership patterns and project future infrastructure needs.